\vspace*{-1cm}
\chapter*{ABSTRACT}
\thispagestyle{fancy}

\begin{center}
\setstretch{1.4}\large \thesistitleA\\ \thesistitleB
\end{center}
\vspace{0.2cm}

\begingroup
\setstretch{1.5}

Crop disaster insurance is difficult to supply privately because weather losses are correlated across farms, so South Korea supports it through premium subsidies, administrative cost subsidies, and state reinsurance. In this thesis, I assess each instrument against its economic function using program rules and financing data. I then compare the cultivated area, yield, and production of seven annual crops with those of crops not yet insured, and rice planted area in early and later pilot municipalities. Proportional premium support ties transfers to insured value and rated risk rather than household income, loading support finances delivery without by itself rewarding efficiency, and state reinsurance shares ordinary as well as catastrophic losses. The mean production contrast for the seven crops is 0.020 log points (95 percent interval [\textminus 0.277, 0.317]); its sign depends on comparison choices, and pre-pilot divergence limits causal interpretation. The rice estimate is \textminus 0.021 log points (exploratory 95 percent interval [\textminus 0.064, 0.022]). The production and planted-area estimates are imprecise and sensitive to comparison choices. The institutional assessment identifies distinct evaluation priorities for the three instruments: access and distribution, delivery cost and quality, and risk allocation and fiscal exposure.

\vspace{0.3cm}
\noindent\textbf{Keywords}: crop disaster insurance, government intervention, premium subsidies, state reinsurance, cultivated area, difference-in-differences, South Korea
\endgroup
